\chapter*{Preface}
A series of books is a city: each tome is a district with its own character, and the streets between them are the
theorems that one book proves and another one uses. This volume is the map of that city.

It is not written like the other tomes. Most of it is generated directly from the sources of the books: every
chapter of every tome, every named theorem with the number it carries in its own tome, and every reference from
one tome to a result of another. At the time of this printing it covers \statTomes\ tomes, \statChapters\
chapters, \statResults\ numbered results, of which \statNamed\ carry a name, \statExercises\ exercises, and
\statCitations\ references between tomes, each one checked to point to a result that exists and has the kind the
reference says.

The first chapter describes the shape of the series: its three parts, the computing track, the companion volumes
and the problem books, what each tome is about and which ones it leans on. The second collects the notation and
the conventions shared by the tomes, and the places where they differ. Then come six chapters, one for each group
of tomes, listing chapter by chapter the named results with their numbers: the place to look when you remember
Dirichlet's theorem or Perron's theorem on positive matrices and want to know where it is proved. The next chapter follows
the threads between the tomes: which tomes cite which, and the results that the rest of the series relies on most.
The last is an alphabetical index of every named result.

Use it as you would use the map of a city you are learning: to find your way back to a theorem you have met, to see
which roads lead to a book you want to read, and to notice how the same few ideas, the prime number theorem, the
central limit theorem, the spectral theorem, Galois theory, connect districts that look far apart.
